\documentclass[11pt,a4paper]{article}
\usepackage{amsmath,amssymb}
\usepackage{siunitx}
\usepackage[margin=2.5cm]{geometry}
\usepackage{graphicx}
\usepackage{booktabs}
\usepackage{xcolor}

\begin{document}

\noindent\textbf{\large Carrier Dynamics Under GHz-Burst Femtosecond Excitation: Piecewise Analytical Solution}\\[4pt]
\noindent Burak Arayici \hfill \today

\bigskip

%=============================================================================
\section{Governing Equation}
%=============================================================================

The free-carrier density inside a silicon nanoparticle under burst excitation evolves as:
\begin{equation}
\frac{dN_\text{fc}}{dt} = \frac{\beta}{2\hbar\omega}\,I^2(t) - \gamma\,N_\text{fc}^3
\label{eq:governing}
\end{equation}
where $\beta$ is the two-photon absorption (TPA) coefficient, $\hbar\omega$ the photon energy, $\gamma$ the Auger recombination coefficient, and $I(t)$ is the instantaneous intensity of the burst pulse train:
\begin{equation}
I(t) = I_0 \sum_{n=0}^{N-1} \exp\!\left[-\frac{4\ln 2}{\tau_\text{p}^2}\,(t - n\,\tau_\text{r})^2\right]
\label{eq:pulse_train}
\end{equation}
Here $I_0$ is the peak intensity, $\tau_\text{p}$ is the pulse FWHM, $\tau_\text{r} = \tau_\text{burst}/(N-1)$ is the inter-pulse spacing, and $N$ is the number of pulses in the burst.

Equation~\eqref{eq:governing} has no general closed-form solution due to the nonlinear $N_\text{fc}^3$ term coupled to the time-dependent source.

%=============================================================================
\section{Timescale Separation}
%=============================================================================

Since $\tau_\text{p} = \SI{100}{fs}$, the Auger term $\gamma N_\text{fc}^3$ is negligible during each pulse --- recombination simply cannot act on such a short timescale. Conversely, between pulses ($I = 0$), only recombination acts. This allows us to solve each phase independently.

%=============================================================================
\section{Phase 1: Carrier Generation During a Pulse}
%=============================================================================

During the $n$-th pulse, we neglect the Auger term:
\begin{equation}
\frac{dN_\text{fc}}{dt} \approx \frac{\beta}{2\hbar\omega}\,I^2(t)
\label{eq:generation}
\end{equation}
Integrating over the full pulse duration (effectively $-\infty$ to $+\infty$ since the Gaussian decays rapidly):
\begin{equation}
\Delta N_\text{fc} = \frac{\beta}{2\hbar\omega} \int_{-\infty}^{\infty} I_0^2 \exp\!\left[-\frac{8\ln 2}{\tau_\text{p}^2}\,t^2\right] dt
\label{eq:deltaN_integral}
\end{equation}
Using the standard Gaussian integral $\int_{-\infty}^{\infty} e^{-at^2}\,dt = \sqrt{\pi/a}$ with $a = 8\ln 2 / \tau_\text{p}^2$:
\begin{equation}
\boxed{\Delta N_\text{fc} = \frac{\beta\,I_0^2\,\tau_\text{p}}{2\hbar\omega}\,\sqrt{\frac{\pi}{8\ln 2}}}
\label{eq:deltaN}
\end{equation}
The factor $\sqrt{\pi/(8\ln 2)} \approx 0.753$ is the Gaussian shape correction relative to a rectangular pulse of the same FWHM and peak intensity.

The carrier density immediately after the $n$-th pulse is:
\begin{equation}
N_n^{+} = N_n + \Delta N_\text{fc}
\label{eq:after_pulse}
\end{equation}
where $N_n$ is the density just before pulse $n$ arrives.

%=============================================================================
\section{Phase 2: Auger Decay Between Pulses}
%=============================================================================

Between pulses, $I = 0$ and Eq.~\eqref{eq:governing} reduces to:
\begin{equation}
\frac{dN_\text{fc}}{dt} = -\gamma\,N_\text{fc}^3
\label{eq:auger_decay}
\end{equation}
This is separable. Separating variables:
\[
\frac{dN_\text{fc}}{N_\text{fc}^3} = -\gamma\,dt
\]
Integrating both sides with initial condition $N_\text{fc}(0) = N_0$:
\[
-\frac{1}{2N_\text{fc}^2}\bigg|_{N_0}^{N(t)} = -\gamma\,t
\quad\Longrightarrow\quad
\frac{1}{N^2(t)} - \frac{1}{N_0^2} = 2\gamma\,t
\]
\begin{equation}
\boxed{N_\text{fc}(t) = \frac{N_0}{\sqrt{1 + 2\gamma\,N_0^2\,t}}}
\label{eq:auger_solution}
\end{equation}

After a time $\tau_\text{r}$ (the inter-pulse gap), the carrier density just before the next pulse is:
\begin{equation}
N_{n+1} = \frac{N_n^{+}}{\sqrt{1 + 2\gamma\,(N_n^{+})^2\,\tau_\text{r}}}
\label{eq:decay_map}
\end{equation}

%=============================================================================
\section{Recurrence Relation}
%=============================================================================

Combining Eqs.~\eqref{eq:after_pulse} and~\eqref{eq:decay_map}, the carrier density evolves pulse-to-pulse as:
\begin{equation}
\boxed{N_{n+1} = \frac{N_n + \Delta N_\text{fc}}{\sqrt{1 + 2\gamma\,\left(N_n + \Delta N_\text{fc}\right)^2 \tau_\text{r}}}}
\label{eq:recurrence}
\end{equation}
with $N_0 = 0$ (no carriers before the first pulse). This replaces the linear accumulation $N_k = k\,\Delta N_\text{fc}$ of the no-recombination model.

%=============================================================================
\section{Steady-State Carrier Density}
%=============================================================================

At steady state, $N_{n+1} = N_n \equiv N_\text{ss}$. Setting Eq.~\eqref{eq:recurrence} equal on both sides gives the implicit equation:
\begin{equation}
\boxed{N_\text{ss} = \frac{N_\text{ss} + \Delta N_\text{fc}}{\sqrt{1 + 2\gamma\,(N_\text{ss} + \Delta N_\text{fc})^2\,\tau_\text{r}}}}
\label{eq:ss_implicit}
\end{equation}
Squaring both sides:
\begin{equation}
N_\text{ss}^2\left[1 + 2\gamma\,\tau_\text{r}\,(N_\text{ss} + \Delta N_\text{fc})^2\right] = (N_\text{ss} + \Delta N_\text{fc})^2
\label{eq:ss_squared}
\end{equation}
which is a quartic in $N_\text{ss}$ --- exact but not amenable to a compact closed-form solution.

\paragraph{Approximate scaling.} In the limit $N_\text{ss} \gg \Delta N_\text{fc}$ (strong Auger clamping), we can approximate $N_\text{ss} + \Delta N_\text{fc} \approx N_\text{ss}$ in Eq.~\eqref{eq:ss_squared}:
\[
N_\text{ss}^2\left(1 + 2\gamma\,\tau_\text{r}\,N_\text{ss}^2\right) \approx N_\text{ss}^2
\]
The left-hand side also contains the $(N_\text{ss} + \Delta N_\text{fc})^2 - N_\text{ss}^2$ difference on the right. Expanding the exact form instead: using the identity $(N_\text{ss} + \Delta N_\text{fc})^2 - N_\text{ss}^2 = \Delta N_\text{fc}(2N_\text{ss} + \Delta N_\text{fc})$ and applying $\Delta N_\text{fc} \ll N_\text{ss}$:
\[
\Delta N_\text{fc}\cdot 2N_\text{ss} \approx 2\gamma\,\tau_\text{r}\,N_\text{ss}^2 \cdot N_\text{ss}^2
\]
\[
\Delta N_\text{fc} \approx \gamma\,\tau_\text{r}\,N_\text{ss}^3
\]
\begin{equation}
\boxed{N_\text{ss} \approx \left(\frac{\Delta N_\text{fc}}{\gamma\,\tau_\text{r}}\right)^{1/3}}
\label{eq:ss_strong}
\end{equation}
This reveals the scaling: steady-state carrier density grows as the cube root of the per-pulse generation rate and decreases with stronger Auger recombination or longer inter-pulse gaps.

%=============================================================================
\section{Per-Pulse Heating}
%=============================================================================

The volumetric heating deposited during the $n$-th pulse has two contributions. TPA directly deposits energy by integrating $\beta I^2$ over the Gaussian envelope:
\begin{equation}
H_\text{TPA} = \beta \int_{-\infty}^{\infty} I^2(t)\,dt = \beta\,I_0^2\,\tau_\text{p}\,\sqrt{\frac{\pi}{8\ln 2}}
\label{eq:H_tpa}
\end{equation}
This is the same for every pulse. Free carriers present before pulse $n$ absorb linearly via $\sigma_\text{fc}\,N_n\,I(t)$, giving:
\begin{equation}
H_{\text{FCA},n} = \sigma_\text{fc}\,N_n \int_{-\infty}^{\infty} I(t)\,dt = \sigma_\text{fc}\,N_n\,I_0\,\tau_\text{p}\,\sqrt{\frac{\pi}{4\ln 2}}
\label{eq:H_fca}
\end{equation}
where $N_n$ is approximately constant during the pulse (generation changes $N_\text{fc}$ by $\Delta N_\text{fc}$, small compared to the accumulated density). The total heating from the $n$-th pulse is therefore:
\begin{equation}
H_n = \beta\,I_0^2\,\tau_\text{p}\,\sqrt{\frac{\pi}{8\ln 2}} + \sigma_\text{fc}\,N_n\,I_0\,\tau_\text{p}\,\sqrt{\frac{\pi}{4\ln 2}}
\label{eq:H_n}
\end{equation}

%=============================================================================
\section{Total Burst Heating and Enhancement}
%=============================================================================

Summing over all $N$ pulses in the burst:
\begin{equation}
H_\text{burst} = N\,\beta\,I_0^2\,\tau_\text{p}\,\sqrt{\frac{\pi}{8\ln 2}} + \sigma_\text{fc}\,I_0\,\tau_\text{p}\,\sqrt{\frac{\pi}{4\ln 2}} \sum_{n=0}^{N-1} N_n
\label{eq:H_burst}
\end{equation}
The first term is the total TPA heating (identical for each pulse), and the second is the cumulative FCA contribution, which depends on the carrier density history $\{N_n\}$ given by the recurrence relation (Eq.~\ref{eq:recurrence}).

The enhancement factor relative to a single nanosecond reference pulse ($H_\text{ns} = \beta\,I_\text{ns}^2\,\tau_\text{ns}$) separates into TPA and FCA contributions:
\begin{equation}
\varepsilon = \varepsilon_\text{TPA} + \varepsilon_\text{FCA}
\label{eq:enhancement}
\end{equation}
The TPA enhancement follows directly:
\begin{equation}
\varepsilon_\text{TPA} = \frac{N\,\beta\,I_p^2\,\tau_\text{p}\,\sqrt{\dfrac{\pi}{8\ln 2}}}{\beta\,I_\text{ns}^2\,\tau_\text{ns}} = N\,\sqrt{\frac{\tau_\text{ns}}{\tau_\text{p}}}\,\sqrt{\frac{\pi}{8\ln 2}}
\label{eq:eps_tpa}
\end{equation}
where the second equality uses $I_p = I_\text{ns}\,(\tau_\text{ns}/\tau_\text{p})^\omega$ with $\omega = 0.75$.

Once the carrier density reaches steady state ($\sum N_n \approx N\,N_\text{ss}$), the ratio of FCA to TPA enhancement reduces to:
\[
\frac{\varepsilon_\text{FCA}}{\varepsilon_\text{TPA}} = \frac{\sigma_\text{fc}\,N_\text{ss}\,I_p\,\tau_\text{p}\,\sqrt{\dfrac{\pi}{4\ln 2}}}{\beta\,I_p^2\,\tau_\text{p}\,\sqrt{\dfrac{\pi}{8\ln 2}}} = \frac{\sqrt{2}\;\sigma_\text{fc}\,N_\text{ss}}{\beta\,I_p}
\]
so the total enhancement factors compactly as:
\begin{equation}
\boxed{\varepsilon = \varepsilon_\text{TPA}\!\left(1 + \frac{\sqrt{2}\;\sigma_\text{fc}\,N_\text{ss}}{\beta\,I_p}\right)}
\label{eq:eps_total}
\end{equation}
where $N_\text{ss}$ is obtained from the implicit steady-state condition (Eq.~\ref{eq:ss_implicit}). The dimensionless ratio $\sigma_\text{fc}\,N_\text{ss}/(\beta\,I_p)$ is the steady-state FCA-to-TPA absorption ratio --- a single number quantifying how much free-carrier effects add on top of the direct TPA enhancement.

\end{document}
